\subsection{The continuous arithmetic kernel and its trace maps}

\begin{statement}{E.33}{Proved}{The continuous arithmetic kernels are free, with explicit generators as a Galois module}
Fix a prime $\ell$ and put $K_m=\mathbb Q(\mu_m)$,
$d_m=[K_m:\mathbb Q]$ and $r_m=r_2(K_m)+1$.
The continuous unipotent completion occurring in E.30 is
\[
 (G_{K_m})^{\rm un}_{\mathbb Q_\ell}
 \simeq \exp\widehat{\mathbb L}(V_m),\qquad
 V_m=\frac{K_m\otimes_{\mathbb Q}\mathbb Q_\ell}
 {\log_\ell(\mathcal O_{K_m}^{\times}\otimes_{\mathbb Z}\mathbb Q_\ell)}.
                                                        \tag{E.33.1}
\]
Here $\widehat{\mathbb L}$ is completion of the free Lie algebra by
bracket length. The vector space $V_m$ is the canonical abelianization
of this Lie algebra; choosing lifts of a basis gives the displayed
noncanonical free presentation. It has dimension $r_m$.
For $m\geq3$, complex conjugation $c$ gives the explicit
$\operatorname{Gal}(K_m/\mathbb Q)$-equivariant identification
\[
 V_m\xrightarrow{\sim}\mathbb Q_\ell\oplus
        (K_m\otimes\mathbb Q_\ell)^-,\qquad
 [z]\longmapsto
 \left(\frac{\operatorname{Tr}_{K_m/\mathbb Q}(z)}{d_m},
                         \frac{z-cz}{2}\right).           \tag{E.33.2}
\]
For $m=1,2$ it is $V_m=\mathbb Q_\ell$.
For $m\mid n$, the actual homomorphism of completed groups is induced
by $G_{K_n}\hookrightarrow G_{K_m}$, as in E.30, and its map on
abelianizations is
\[
 [z]\longmapsto[\operatorname{Tr}_{K_n/K_m}z].             \tag{E.33.3}
\]
In (E.33.2) this is $(a,w)\mapsto([K_n:K_m]a,
\operatorname{Tr}_{K_n/K_m}w)$. Thus every transition is faithfully
flat. These groups and these specified homomorphisms compute the
continuous arithmetic kernel in (E.30.2); no identification of
the abstract rational completion with this continuous completion is
used.
More explicitly, along the cofinal conductors $m_j=4j!$, free
generators can be chosen compatibly so that each transition keeps
the first $r_{m_j}$ generators and sends every additional generator
to $1$. Thus that kernel is an inverse limit of free prounipotent
groups of the explicit ranks $r_{m_j}$ with these deletion maps.
\end{statement}
\begin{proof}
Write $K=K_m$, and let $S$ consist of the places above $\ell$ and
the infinite places. Every continuous homomorphism
$\rho:G_K\to U(\mathbb Q_\ell)$, for $U$ unipotent algebraic,
is unramified outside $S$. Indeed its compact image is pro-$\ell$:
use a central series of $U$ with additive quotients and the fact
that a compact additive subgroup of a finite-dimensional
$\mathbb Q_\ell$-space is pro-$\ell$. At a finite place $v\nmid\ell$,
wild inertia therefore dies, and tame inertia maps through its
$\mathbb Z_\ell(1)$ quotient. If $N$ is the logarithm of the image
of a tame generator and $F_v$ is arithmetic Frobenius, the exact
relation $F_v\tau F_v^{-1}=\tau^{q_v}$ gives
\[
 \operatorname{Ad}(\rho(F_v))N=q_v N.
\]
The left-hand operator is unipotent and $q_v\ne1$ in
$\mathbb Q_\ell$. Hence $N=0$. The normal closed subgroup generated
by these inertia groups is the kernel of $G_K\to G_{K,S}$,
which proves the claimed factorization, with its universal property.

The class-field-theoretic exact sequence of local principal units,
global units and the maximal abelian pro-$\ell$ extension unramified
outside $S$ gives
\[
 H_1(G_{K,S},\mathbb Q_\ell)
 =\operatorname{coker}\left(
 \mathcal O_K^{\times}\otimes\mathbb Q_\ell
 \xrightarrow{\log_\ell}
 \bigoplus_{v\mid\ell}K_v\right).                       \tag{E.33.4}
\]
The finite ideal class group disappears after tensoring. The exact
sequence, including its proof from reciprocity, is
\cite[Theorem1.5.4 and Theorem1.5.7]{E_SharifiIwasawa}.
The diagonal logarithm is injective here: $K_m/\mathbb Q$ is
abelian, so the proved abelian-field case of Leopoldt applies
\cite[Theorems1.5.17 and1.5.21]{E_SharifiIwasawa}.
To specify the reduction, the unit representation of the maximal
real subfield is its augmentation representation. A rational
normal generator gives conjugate units whose logarithms have just
the norm relation over $\mathbb Q$; Brumer's linear-independence
theorem retains independence over algebraic coefficients.
Diagonalizing the abelian group matrix by characters makes every
nontrivial character eigenvalue nonzero. Thus the logarithm has
the full Dirichlet rank. This is precisely the proof of the cited
Theorem1.5.21, and applies to every $\ell$, including $2$.
Consequently (E.33.4) has dimension
$[K:\mathbb Q]-(r_1(K)+r_2(K)-1)=r_2(K)+1$.

We also need degree two, rather than infer freeness from this
dimension. Apply the finite-coefficient Euler characteristic
\cite[TheoremI.5.1]{E_MilneADT} to $\mathbb Z/\ell^a$ with
trivial action. Its formula is exactly
\[
 \frac{|H^0(G_{K,S},\mathbb Z/\ell^a)|
       |H^2(G_{K,S},\mathbb Z/\ell^a)|}
      {|H^1(G_{K,S},\mathbb Z/\ell^a)|}
       =\ell^{-a r_2(K)}.                               \tag{E.33.5}
\]
All these finite-coefficient groups are finite, by
\cite[CorollaryI.4.15]{E_MilneADT}. The class-field computation says
that $H^1$ has order $\ell^{a(r_2+1)}$ times a bounded factor,
because its abelian pro-$\ell$ source is finitely generated with
finite torsion. Since $|H^0|=\ell^a$, (E.33.5) bounds the orders
of $H^2$ independently of $a$. Inverse limits are exact on these
finite systems; the resulting $H^2(G_{K,S},\mathbb Z_\ell)$
is torsion, and tensoring with $\mathbb Q_\ell$ gives
\[
 H^2_{\rm cts}(G_{K,S},\mathbb Q_\ell)=0.                \tag{E.33.6}
\]
Continuous cochains with $\mathbb Q_\ell$ coefficients are the
union of the cochain complexes with values in
$\ell^{-b}\mathbb Z_\ell$, by compactness of their domain; this
justifies the last tensoring step. The real places at $\ell=2$
cause no exception: (E.33.5) is exactly Milne's degree-zero-through-two
formula, whose footnote explicitly treats this case.

For completeness, the universal unipotent completion has the same
$H^1$ and injects on $H^2$ into continuous group cohomology.
The first assertion follows from homomorphisms to the additive
group. For the second, an algebraic continuous Lie-cohomology
class factors through a finite-dimensional quotient. Represent it
by a central extension by $\mathbb G_a$; exponential coordinates
give a continuous set-theoretic section of the projection.
Its pullback along the dense continuous
group homomorphism is a continuous group extension. A zero pullback
class supplies a continuous lift; the universal property of
unipotent completion extends that lift algebraically. Its composite
with the projection is the identity on the dense original group,
hence the identity, so it splits the original extension. This
proves injectivity, with no comparison theorem left implicit.
Equation (E.33.6) therefore makes the second cohomology of the
completed Lie algebra zero.

Choose lifts of a basis of its abelianization $V_m$. They give a
surjection from $\widehat{\mathbb L}(V_m)$, with a closed kernel
$R$ contained in its commutator ideal. The free-Lie five-term
sequence identifies the continuous dual of
$R/\overline{[\widehat{\mathbb L}(V_m),R]}$ with degree-two cohomology.
A nonzero $R$ has a nonzero term of least bracket degree, which
cannot lie in $\overline{[\widehat{\mathbb L}(V_m),R]}$. Thus $R=0$,
proving (E.33.1).

For the more explicit module in (E.33.2), a unit in a CM field has
unit/conjugate-unit a root of unity. Its logarithm is therefore
fixed by $c$ and has trace zero. The injectivity already proved
and the equal dimensions $d_m/2-1$ show that the logarithmic image
is exactly the trace-zero subspace of
$(K_m\otimes\mathbb Q_\ell)^+$. The kernel and image of the
displayed formula are now as asserted. The scalar $1/2$ is
legitimate also over $\mathbb Q_2$.

Finally reciprocity sends the inclusion $G_{K_n}\to G_{K_m}$
to the local norm maps. Logarithms send norms to traces and send
norms of global units to the logarithms of global units. This
proves (E.33.3), including its compatibility with every tower.
Trace is onto on the displayed quotient: on an embedded element
of $K_m$ it is multiplication by the nonzero degree. A morphism
of prounipotent groups surjective on this finite-dimensional
abelianization is surjective, successively on the lower central
quotients. The transition maps are consequently faithfully flat.
To obtain the final presentation, choose generators at $m_1$.
At the next stage lift their logarithms through the surjection.
Choose a basis of the remaining kernel on abelianizations and
lift it into the Lie kernel: an arbitrary lift can be corrected
by a commutator, because the map on commutator ideals is onto.
Together these elements form a basis modulo commutators, and hence
a free set of topological Lie generators by the freeness already
proved. Iterating gives exactly the claimed deletion maps.

The accompanying exact cyclotomic-field program verifies the
conjugation projections, all divisor trace maps and their tower
compositions through conductor $24$, and prints the corresponding
free-Lie ranks through degree $8$. This checks the explicit maps,
not merely their dimensions.
\end{proof}
